\section{Introduction}\label{intro}

Given a (positive Borel) measure $\mu$ on $\R$, a classical problem in numerical analysis is to approximate the integral with respect to $\mu$ of a suitably well-behaved function~$f$. One approach is via so called \emph{quadrature rules}. These approximate the integral by a weighted sum of function values at specified points. One classical construction for quadrature rules designed to approximate the integral of continuous functions consists of demanding an exact evaluation of the integral for all polynomials of degree $\leq D$. If the moments of $\mu$ exist and are finite, then this amounts to finding a measure supported on finitely-many points whose moments agree with those of $\mu$ up to degree~$D$.

We use $\t$ as a formal variable on the real line and write $\R[\t]_{\leq\,D}$ for the vector space of real polynomials of degree at most $D$. For $k \in
\N_0$, we denote the $k^{\rm th}$ moment of $\mu$, if it exists and is finite, by
\[
m_k \quad = \quad \int \t^k \ d\mu.
\]

\begin{defi}\label{dfqr}
Suppose $D$ is a positive integer and $\mu$ is a measure on $\R$ whose moments up to degree~$D$ exist and are finite. For $x \in \R \cup \{\infty\}$, define the linear function
\[
\ev_x:\R[\t]_{\leq D} \rightarrow \R
\]
as follows. For $f = \sum_{k\,=\,0}^Df_k \t^k \in \R[\t]_{\leq\,D}$ with $f_0,\,\ldots,\,f_D\in\R,$
\[
\ev_x(f) = f(x) \quad \text{ for }x\in \R \quad \text{and } \quad \ev_{\infty}(f) = f_D.
\]
We sometimes write $\ev_{\infty}^D$ for $\ev_{\infty}$ to emphasize its dependence on $D$. A \emph{quadrature rule} of degree~$D$ for $\mu$ is a finite set $N\subset \R\cup\{\infty\}$ together with a function $w\colon N\to\R_{>0}$ with
\[
\int f \; d \mu \, = \, \sum_{x\in N}w(x) \ev_x(f) \quad \text{ for all }\quad f\in\R[\t]_{\leq\,D}.
\]
We call $N$ the \emph{nodes} of the quadrature rule.
\end{defi}


\begin{rema}
In Definition~\ref{dfqr} we allow nodes at infinity. This is not desirable for application in numerical analysis, since one cannot generally evaluate functions at these nodes. However, compactifying the real line by adding $\infty$ makes certain arguments easier. Although our statements are phrased with this more general notion of quadrature, our main result presented below actually provides a tool to explicitly distinguish the cases of quadrature rules with all nodes real from the case where one of the nodes is $\infty$. Furthermore, it is a classical theorem that quadrature rules for $\mu$ with no nodes at infinity exist for every degree~$D$ provided the moments of $\mu$ exist and are finite up to the same degree (see e.g.~\cite[Theorem~5.8]{la1}, \cite[Theorem~5.9]{la2}, \cite[Theorem~1.24]{sch}).
\end{rema}

We say the measure $\mu$ is \emph{non-degenerate} in degree $d$ if its moments $m_k$ are finite up to degree $k=2d$ and for every nonzero, nonnegative polynomial $ f\in \R[\t]_{\leq\,2d}$ we have $\int f \ d\mu >0$. Since in one variable any nonnegative polynomial $f$ is sum of squares of polynomials, this is equivalent to demanding that $\int p^2 \ d\mu >0$, for every $0\neq p\in \R[\t]_{\leq\,d}$. This property can be checked quite conveniently in the following~way.


\begin{defi} \label{def:Mmatrix}
Consider the quadratic form $p\in \R[\t]\mapsto\int p^2 d\mu$ and restrict it to $\R[\t]_{\leq\,d}$. With respect to the monomial basis $1,\t,\hdots, \t^d$ for $\R[\t]_{\leq\,d}$, this quadratic form is represented by the $(d+1)\times (d+1)$ Hankel matrix with $(i,j)$th entry $m_{i+j-2}$:
\[
M_d \; = \;
\begin{pmatrix}
m_0 & m_1 & m_2 & \hdots &m_d \\
m_1 & m_2 & \iddots&&m_{d+1} \\
m_2 & \iddots& &\iddots&\vdots\\
\vdots&&\iddots&&m_{2d-1}\\
m_d &m_{d+1}&\hdots&m_{2d-1}&m_{2d}
\end{pmatrix}.
\]
With this notation, a measure possessing finite moments up to degree~$2d$ is non-degenerate in degree $d$ if and only if $\det(M_d)\neq 0$.
\end{defi}


From the point of view of numerical analysis it is desirable to have a quadrature rule that is exact up to a certain degree and requires the fewest number of evaluations possible. Such a quadrature rule is called a \emph{Gaussian quadrature rule} for $\mu$.

It is known that when $D=2d+1$ is odd and $\mu$ is non-degenerate in degree $d$, there is a unique quadrature rule for $\mu$ with $d+1$ nodes, and none with fewer nodes~\cite[Chapter~I, Theorem.~3]{anheizer}, see also~\cite[Theorem.~3.8]{curto}. The nodes can be found as follows. Let $M'_d$ denote the $(d+1)\times(d+1)$ matrix representing the quadratic form $p\mapsto \int \t \cdot p^2 d \mu$ with respect to the monomial basis $1,\t,\hdots, \t^d$ of $\R[\t]_{\leq \,d}$, meaning that $(M'_d)_{i,j}$ equals $m_{i+j-1}$. Then the $d+1$ nodes of the unique Gaussian quadrature rule for $\mu$ in degree~$2d+1$ are the $d+1$ roots of the univariate polynomial $\det(\x M_d - M_d') $, see e.g.~\cite[formule 2.2.9, p.~27]{sze}. Since $M_d$ and $M'_d$ are real symmetric matrices and $M_d$ is positive definite, this writes the problem of finding these $d+1$ nodes as a
\emph{generalized eigenvalue problem}.


In this paper we focus on the case when $D=2d$ is even. In this case, there is a one-parameter family of quadrature rules for $\mu$ with $d+1$ nodes (see e.g.~\cite[Theorem~9.7]{sch}). Here we reprove this fact by constructing a polynomial $F\in \R[\x,\y]$ with degree $d$ in each of $\x$ and $\y$ and the property that for every $y\in \R$, the $d$ roots of $F(\x,y)\in \R[\x]_{ \leq\,d}$ are the other $d$ nodes (among them possibly $\infty$) of this unique quadrature rule with $y$ as a node.

Furthermore, we give symmetric determinantal representations of $F$, which again translates the problem of finding nodes of a quadrature rule into finding the generalized eigenvalues of a real symmetric matrix,
i.e. solving $\det(\x\,A-B)=0$ where $A,B$ are real symmetric matrices and $A$ is positive semidefinite (see e.g.~\cite[Chapter~4]{eig} or~\cite{golub}). While the literature on Gaussian quadrature rules is vast, to our knowledge these formulas are new. A result in the same spirit was obtained in~\cite[Proposition~3.3]{curto}. However, the procedure advertised in~\cite[Theorem~3.10$\MK$(ii)]{curto} requires to compute the zeros of a univariate polynomial which is more expensive than solving a generalized eigenvalue problem~\cite[p.~1225]{golub}. Moreover, it involves inversion of Hankel matrices which is numerically less stable, see~\cite{illhankel}.


To construct $F$ and its determinantal representations, we consider three quadratic forms on $\R[\t]_{\leq\,d-1}$, namely those taking $p\in \R[\t]_{\leq\,d-1}$ to $\int p^2d\mu$, $\int \t \cdot p^2d\mu$, and $\int \t^2 \cdot p^2d\mu$, respectively. We call $M_{d-1}$, $M'_{d-1}$ and $M''_{d-1}$ the matrices representing these quadratic forms with respect to the basis $1,\t,\hdots, \t^{d-1}$. That is, define the $d\times d$ ~matrices
\begin{equation}\label{eq:matrices}
\begin{aligned}
M_{d-1}\; &= \; (m_{i+j-2})_{1\,\leq\,i,j\, \leq\,d}, \\
M'_{d-1} \; &= \; (m_{i+j-1})_{1\,\leq\,i,j\, \leq\,d}, \; \\ 
\text{ and } M''_{d-1} \,& = \, (m_{i+j})_{1\,\leq\,i,j\,\leq\,d}.
\end{aligned}
\end{equation}

\begin{theo}\label{thm:main}
Let $\mu$ be a measure on $\R$ that is non-degenerate in degree $d\geq1$. There is a unique (up to scaling) polynomial $F\in \R[\x,\y]$ of degree~$2d$ with the property that for $x,y\in \R$, $F(x,y)=0$ if and only if there is a Gaussian quadrature rule for $\mu$ of degree~$2d$ with nodes $x=r_1,y=r_2,r_3, \hdots, r_{d+1} $ in $ \R\cup \{\infty\}$. This polynomial has the following two determinantal representations:
\begin{enumerate}\alphenumi
\item \label{list1.4.1}the determinant of the $d\times d$ matrix with bilinear entries in $\x,\y$,
\[
F \, = \, \det\left(\x\y M_{d-1} - (\x+\y)M'_{d-1} + M''_{d-1}\right),
\]
\item \label{list1.4.2}the determinant of the $(2d+1)\times (2d+1)$ matrix with linear entries in $\x,\y$,
\[
(\x-\y) \cdot F \, = \, c\cdot \det
\begin{pmatrix}
\frac{\det M_{d-1}}{\det M_d} (\x-\y) &e_d^T&e_d^T\\
e_d&\x M_{d-1}-M_{d-1}'&0\\
e_d&0&-\y M_{d-1}+M_{d-1}'
\end{pmatrix}
\]
where $e_d=(0,\ldots,0,1)^T\in\mathbb{R}^d$ and $c=(-1)^d \det(M_d)/\det(M_{d-1})^2$.
\end{enumerate}
Moreover, for all $x\in \R$, with $\det(x M_{d-1} -M_{d-1}')\neq 0$, all nodes are on the real line, i.e., the associated quadrature rule has no nodes at infinity.
\end{theo}

We prove the two parts of this theorem in Sections~\ref{sec:BilinearDetRep} and~\ref{sec:LinearDetRep}. The implications for finding quadrature rules as generalized eigenvalues are made explicit in Remarks~\ref{rem:GenEigBilinear} and~\ref{rem:GenEigLinear}. In Section~\ref{sec:Generalizations}, we discuss possible generalizations of this result.

\subsection*{Acknowledgments}
This collaboration was initiated during the Oberwolfach workshop ``Real algebraic geometry with a view toward moment problems and optimization'' in March 2017. Our thanks go to the organizers of the workshop and to the Mathematisches Forschungsinstitut Oberwolfach. We thank Petter Br\"andén, Christoph Hanselka, Rainer Sinn and Victor Vinnikov for interesting discussions


\section{Bilinear Determinantal Representation}\label{sec:BilinearDetRep}
In this section, we prove Theorem~\ref{thm:main}$\MK$\eqref{list1.4.1}. The proof relies heavily on understanding a particular symmetric bilinear form $B_{x,y}$ on $\R[\t]_{\leq\,d-1}$. For $x,y\in \R$ and\linebreak $p,q\in \R[\t]_{\leq\,d-1}$, let
\[
B_{x,y}(p,q) \ = \ \int (x-\t)(y-\t) p q \ d\mu \ = \ xy\int p q \ d\mu-(x+y)\int \t\cdot p q \ d\mu+\int \t^2\cdot p q \ d\mu.
\]
Note that the symmetric matrix $xy M_{d-1} - (x+y)M'_{d-1} + M''_{d-1}$ represents $B_{x,y}$ with respect to the basis $1,\t,\t^2,\hdots, \t^{d-1}$. Define $F\in \R[\x,\y]$ to be the polynomial
\[
F\ \ = \ \ \det(\x\y M_{d-1} - (\x+\y)M'_{d-1} + M''_{d-1}).
\]
We will show that $F$ satisfies the requirements of Theorem~\ref{thm:main}.

\begin{proof}[Proof of Theorem~\ref{thm:main}$\MK$\eqref{list1.4.1}]
Suppose that there exists a quadrature rule for $\mu$ of degree~$2d$ with nodes $r_1, r_2,\,\hdots,\,r_{d+1}$ in $\R\cup\{\infty\}$, meaning that there exist $w_1,\,\hdots,\, w_{d+1}\in \R_{>\,0}$ so that
\begin{equation}\label{eq:nodes}
\int f \ d\mu \; = \; \sum_{i=1}^{d+1} w_i \ev_{r_i}(f) \; \text{ for all } \; f\in \R[\t]_{\leq\,2d}.
\end{equation}
Note that because $\mu$ is non-degenerate in degree $d$, the nodes $r_1,\,\hdots,\,r_{d+1}$ are necessarily distinct, so that after reindexing we may assume that $r_1,\,\hdots,\, r_d\in \R$.

We claim that $F(r_1, r_2) = 0$. To see this, let $q$ be the unique (up to scaling) nonzero polynomial with $\deg(q)\leq\,d-1$ and $\ev_{r_i}(q)=0$ for each $i=3,\,\hdots,\, d+1$. If each $r_i\in \R$, then we can take $ q$ to be $(\t-r_3)\,\hdots\,(\t-r_{d+1})\in \R[\t]_{\leq\,d-1}$. If $r_{d+1}=\infty$, then we can take $ q= (\t-r_3)\,\hdots\,(\t-r_{d})$. For any $p\in \R[\t]_{\leq\,d-1}$, it follows that
\[
B_{r_1,r_2}(p,q) \; = \; \int (r_1-\t)(r_2-\t) p q \; d\mu \; =\ 0.
\]
The last equality follows from~\eqref{eq:nodes} and the fact that $\deg(p)\leq d-1$. Therefore $q$ is an element of the right kernel of $B_{r_1, r_2}$. Since $B_{r_1, r_2}$ drops rank, the determinant $F(r_1, r_2)$ of the representing matrix equals zero.

Conversely, suppose that for $x,y\in \R$, $F(x,y)=0$. Then the kernel of $B_{x,y}$ contains a polynomial $q\in \R[\t]_{\leq\, d-1}$. For all $p\in \R[\t]_{\leq\,d-1}$,
\[
\int (\t-x)(\t-y)p q\ d\mu = 0.
\]
First, we argue that $x, y,$ and the roots of $q$ are real and pairwise distinct and that the degree of $q$ is $d-2$ or $d-1$. If not, there would exist a non-zero polynomial $p\in \R[\t]_{\leq\,d-1}$ for which $f=(\t-x)(\t-y)p q$ is nonnegative on $\R$. Since $\int f \ d\mu =0$, this contradicts the assumption that $\mu$ is non-degenerate in degree $d$.

Let $r_1=x$, $r_2=y$ and denote the roots of $q$ by $r_3,\,\hdots,\,r_{d+1}$, where we take $r_{d+1}=\infty$ if $\deg(q) = d-2$. Consider the conic hull of the $d+1$ points $\ev_{r_1},\,\hdots,\,\ev_{r_{d+1}}$ in $\R[\t]_{\leq\,2d}^*$. This is a $(d+1)$-dimensional simplicial convex cone in the $(2d+1)$-dimensional vector space $\R[\t]_{\leq\,2d}^*$. Therefore this cone is defined by $d$ linear equalities and $d+1$ linear inequalities, which we will now identify by inspection. For each $i=1, \hdots, d+1$, let $f_i$ be the unique (up to scaling) nonzero polynomial of degree $\leq d$ for which $\ev_{r_j}(f_i) = 0$ for each $j\neq i$. For example, $f_{d+1} =
\prod_{i=1}^{d}(\t-r_i)$. Note that the polynomials $(\t-r_1)(\t-r_2)q\,\cdot\,
\t^k$ for $0\le k\le d-1$ and $f_i^2$ for $1\le i\le d+1$ are linearly independent in $\R[\t]_{\leq\,2d}$. It therefore follows that the conic hull of $\ev_{r_1},\,\hdots,\, \ev_{r_{d+1}}$ is the set of $L\in \R[\t]_{\leq\,2d}^*$ satisfying
\[
L((\t-r_1)(\t-r_2)pq) = 0 \; \text{ for all } \ p \in \R[\t]_{\leq\,d-1}
\quad\text{and} \quad L\left(f_i^2\right) \geq 0 \; \text{for} \; i=1,\,\hdots\,d+1.
\]
The linear function $L_{\mu}\in \R[\t]_{\leq\, 2d}$ given by $L_{\mu}(f) = \int f\, d\mu$ belongs to this convex cone. Hence there are nonnegative weights $w_1,\,\hdots,\,w_{d+1}$ for which
\[
L_{\mu}(f) = \sum_{i=1}^{d+1}w_i \ev_{r_i}(f).
\]
Since $\mu$ is non-degenerate in degree $d$, each of these weights must be positive.
\end{proof}



\begin{rema}\label{rem:PDdiag}
Let $M = M(\x, \y) = \x\y M_{d-1}-(\x+\y)M'_{d-1}+ M''_{d-1}$. We remark that for every $x\in \R$, the matrix $M(x,x)$ is positive definite. To see this, note that $M(x,x)$ represents the quadratic form on $\R[\t]_{\leq\,d-1}$ given by $p \mapsto \int p^2(\t-x)^2 d\mu$ with respect to the monomial basis. Since $\mu$ is non-degenerate in degree $d$, this is positive for any $0\neq p \in \R[\t]_{\leq\,d-1}$.
\end{rema}


\begin{rema}\label{rem:GenEigBilinear}
For fixed $y\in \R$, this allows us to find the roots of $F(\x,y)\in \R[\x]$ as the generalized eigenvalues of a $d\times d$ real symmetric matrix. If $y$ is larger than all of the roots of $\det(\y M_{d-1} - M'_{d-1})$, then $yM_{d-1} - M'_{d-1}$ is positive definite, meaning that $M(\x,y)$ has the form $\x A-B$ where $A$ is a positive definite matrix. We can find the roots in $\x$ as the following generalized eigenvalue problem:
\[
\{x : F(x,y)=0\} \, = \, \left\{ x \, : \, \det\left(x\left(yM_{d-1} - M'_{d-1}\right) - \left(yM'_{d-1} -M_{d-1}''\right)\right)=0\right\}.
\]
If $yM_{d-1} - M'_{d-1}$ is not positive definite, this formula still holds, but may not be as numerically stable. We can instead make a change of variables $\x=\tilde{\x}+y$, which gives
\begin{align*}
M(\x,y) \, = \, M(\tilde{\x}+y,y) &\, = \, (\tilde{\x}+y)yM_{d-1}-(\tilde{\x}+2y)M'_{d-1}+ M''_{d-1} \\
& \, = \, \tilde{\x}\left(yM_{d-1} - M'_{d-1}\right) + M(y,y).
\end{align*}
In particular, this has the form $\tilde{\x}A+B$ where $B$ is positive definite. Suppose $\lambda$ is a root of $\det(\blambda B- A)$. Then $\tilde{\x}=-1/\lambda$ is a solution to $\det(\tilde{\x}A+B)=0$. Therefore
\[
\{x : F(x,y)=0\} \, = \, \left\{ y-1/\lambda \, : \, \det\left(\lambda M(y,y) - yM_{d-1} + M'_{d-1}\right)=0\right\}.
\]
Note that this even works when $\lambda = 0$ if we take $y-1/\lambda$ to be $\infty$.
\end{rema}

\begin{coro}\label{cor:Uniqueness}
Let $\mu$ be a measure on $\R$ that is non-degenerate in degree $d\geq1$. For $y\in \R$ there is a unique quadrature rule for $\mu$ of degree~$2d$ with $d+1$ nodes, one of which is $y$.
\end{coro}

\begin{proof}
Let $y\in \R$. The polynomial $F(\x,y) = \det(M(\x,y)) \in \R[\x]$ has degree at most $d$. Since the matrix pencil $\{M(x,y): x\in \R\}$ contains a positive definite matrix $M(y,y)$, the roots of $F(\x,y)=0$ must all be real. In particular the existence of one real root $x$ implies, by Theorem~\ref{thm:main}, the existence of a quadrature rule of degree~$2d$ of $\mu$ whose $d+1$ nodes include $y$. Moreover the other $d$ nodes $x, r_3,\,\hdots,\,r_{d+1}$ are necessarily the $d$ roots of $F(\x,y)=0$. As in the proof of Theorem~\ref{thm:main}, the conic hull of $\ev_x, \ev_y, \ev_{r_3},\,\hdots,\,\ev_{r_{d+1}}$ is a simplicial cone containing $L_{\mu}$, meaning that there is a unique representation of $L_{\mu}$ as a nonnegative combination of them.
\end{proof}


\begin{exam}
[Normal Distribution, $d=3$] \label{ex:Normal}
Let $\mu$ be the normal Gaussian distribution on $\R$ with mean $0$ and variance $1$. Its moments are given by $m_{2i+1}=0$ and $m_{2i} = (2i-1)!!$ for $i\in \N_0$. For example, the first seven moments of this measure are $m_1=m_3=m_5=0$, $m_0=m_2=1$, $m_4=3$, and $m_6 = 15$. For $d=3$, the polynomial $F$ given by Theorem~\ref{thm:main} is the determinant of the $3\times3$ matrix polynomial
\[
M \, = \,
\begin{pmatrix}
\x \y+1 & -(\x+\y) & \x \y+3 \\
-(\x+\y) & \x \y+3 & -3 (\x+\y) \\
\x \y+3 & -3 (\x+\y) & 3 \x \y+15 \\
\end{pmatrix}.
\]
For fixed $y\in \R$, the polynomial $F(\x,y)\in \R[\x]$ has three real roots, $r_1, r_2, r_3 \in \R\cup\{\infty\}$, which, together with $y$, form the nodes of a quadrature rule for $\mu$ of degree~$6$. The matrix $M(\x,y)$ has the form $\x A+B$ for some real symmetric matrices $A$ and $B$ (depending on $y$). The roots of $F(\x,y)$ are the roots of $\det(\x A + B)$.

\begin{figure}
\begin{center}
%\includegraphics[height=3in]{SexticBL}
\includegraphics[height=3in]{Figures/SexticBL.pdf} 
\end{center}
\caption{\label{fig:Normal}
The sextic curve given by $F=0$, the line $y=1$, and signs of eigenvalues of the $3\times 3$ matrix polynomial $M$ from Example~\ref{ex:Normal}. }
\end{figure}

Moreover, by making a change of coordinates we can make $A$ positive definite. For example, for $y=1$, we set $\x =1/\blambda+1$ to get
\[
\blambda \cdot M(1/\blambda+1,1) =
\begin{pmatrix}
2 \blambda+1 & -2 \blambda-1 & 4 \blambda+1 \\
-2 \blambda-1 & 4 \blambda+1 & -6 \blambda-3 \\
4 \blambda+1 & -6 \blambda-3 & 18 \blambda+3 \\
\end{pmatrix},
\]
which has the form $\blambda A-B$, where $A=M(1,1)$ is positive definite. Solving the generalized eigenvalue problem $\det(\blambda A - B)=0$ gives $\blambda \approx -0.66, -0.32, 0.60$. The solutions to $F(\x,1)=0$ are then $\x = 1+1/\blambda \approx-2.15, -0.52, 2.67$. Thus there is a quadrature rule for $\mu$ of degree~$6$ with nodes consisting of $1$ and these three roots. The curve given by $F(x,y)=0$ along with the line $\y=1$ are shown in Figure~\ref{fig:Normal}.

For $y\in \{0,\pm \sqrt 3\}$, the polynomial $F(\x,y)\in \R[\x]$ becomes quadratic with two real roots $\{0,\pm \sqrt{3}\}\backslash\{y\}$. From this we conclude that there is a quadrature rule of degree~$6$ for $\mu$ with the four nodes $0,\pm \sqrt 3, \infty$.

For $y>\sqrt{3}$, the matrix coefficient of $\x$ in $M(\x,y)$ is positive definite. In this case, $M(\x,y)$ already has the form $\x A-B$ where $A$ is positive definite, so no change of coordinates is required to translate this into a generalized eigenvalue problem.
%\hfill $\diamond$
\end{exam}



\section{Linear Determinantal Representation}\label{sec:LinearDetRep}

In this section we prove Theorem~\ref{thm:main}$\MK$\eqref{list1.4.2}. As in Section~\ref{sec:BilinearDetRep}, we construct a bilinear form depending on $x,y\in \R$ and construct a non-zero kernel of it for those pairs $x,y$ that can be extended to $d+1$ nodes of a quadrature rule for $\mu$ of degree~$2d$.

For $x,y\in \R$, we define a bilinear form $B_{x,y}$ on $\R \oplus \R[\t]_{\leq\,d-1} \oplus \R[\t]_{\leq\,d-1} \cong \R^{2d+1}$. Given $p = (p_0, p_1, p_2)$ and $q = (q_0,q_1,q_2)$ in $\R \oplus \R[\t]_{\leq d-1} \oplus \R[\t]_{\leq\,d-1}$, define
\begin{multline*}
B_{x,y}(p,q) = \int p_1q_1 (x-\t) + p_2q_2 (\t-y) d\mu \\
+ \ev^{d-1}_{\infty}(q_0(p_1+p_2)+p_0(q_1+q_2)) +\frac{p_0q_0\det(M_{d-1})}{\det(M_d)}(x-y).
\end{multline*}
Choosing the basis $1,\t,\,\hdots,\,\t^{d-1}$ for both copies of $\R[\t]_{\leq\,d-1}$ represents this symmetric bilinear form as the $(2d+1)\times (2d+1)$ matrix given in Theorem~\ref{thm:main}$\MK$\eqref{list1.4.2}. That is, if $\vec f$ denotes the coefficients of $f\in \R[\t]_{\leq\,d -1}$ so that $f = \vec{f} \cdot (1,\t,\hdots, \t^{d-1})$, then
\[
B_{x,y}(p,q) \, =
\begin{pmatrix}
p_0 \\ \vec{p_1} \\ \vec{p_2}
\end{pmatrix}^T
\begin{pmatrix}
\frac{\det M_{d-1}}{\det M_d} (x-y) &e_d^T&e_d^T\\
e_d&xM_{d-1}-M_{d-1}'&0\\
e_d&0&-yM_{d-1}+M_{d-1}'
\end{pmatrix}
\begin{pmatrix}
q_0 \\ \vec{q_1} \\ \vec{q_2}
\end{pmatrix},
\]
where $M_{d-1}$ and $M'_{d-1}$ are the $d\times d$ matrices defined in~\eqref{eq:matrices}.


In order to construct a non-zero element in the kernel of $B_{x,y}$, we need to build up some basic facts. The first concerns quadrature rules for $\mu$ whose nodes include $\infty$. As this polynomial will be used heavily in the text below, denote
\[
F_{\infty} = \det(\x M_{d-1} -M_{d-1}').
\]


\begin{lemm}\label{lem:InftyNodes}
The polynomial $F_{\infty}$ has $d$ real roots $s_1, \hdots, s_d \in \R$ and there exist weights $w_1, \hdots, w_d\in \R_{> 0}$ for which
\begin{equation}\label{eq:InftyQuad}
\int f \ d\mu \, = \, w_{\infty} \ev_{\infty}^{2d}(f) + \sum_{i=1}^d w_i \ev_{s_i}(f) \, \text{ for all } \, f\in \R[\t]_{\leq\,2d},
\end{equation}
where $w_{\infty} = \det(M_d)/\det(M_{d-1})$.
\end{lemm}


\begin{proof}
Since $M_{d-1}$ is positive definite, $F_{\infty}$ has $d$ real roots $s_1, \hdots, s_d\in \R$, up to multiplicity. For any root $s$, the matrix $sM_{d-1} - M'_{d-1}$ has rank $<d$, meaning that the polynomial
\[
F(\x,s) = \det(\x(sM_{d-1} - M'_{d-1}) - (sM'_{d-1} -M_{d-1}''))
\]
has degree $\leq d-1$ in $\x$. For any $r\in \R$ with $F(r,s)=0$, there is a quadrature rule for $\mu$ of degree~$2d$ with $d+1$ nodes containing $s$, $r$, and $\infty$. Then the unique quadrature rule of degree~$2d$ with $d+1$ nodes containing $r$ also contains the node $\infty$. This implies that $F(\x,r)$ has degree $\leq d-1$ and $F_\infty(r)=0$, meaning that $r\in \{s_1, \hdots, s_d\}$. Therefore there is a quadrature rule of degree~$2d$ for $\mu$ with nodes $s_1, \hdots, s_d, \infty$. In particular, there exist $w_1, \hdots, w_d, w_\infty\in \R_{>0}$ for which~\eqref{eq:InftyQuad} holds.

Then $w_\infty$ is the largest $\lambda$ for which the quadratic form $p \mapsto \int p^2 d\mu - \lambda \ev_{\infty}(p^2)$ is nonnegative on $\R[\t]_{\leq d}$. This is the largest $\lambda$ for which the matrix $M_d - \lambda e_{d+1}e_{d+1}^T$ is positive semidefinite. We find this by solving the equation $\det(M_d - \blambda e_{d+1}e_{d+1}^T) = 0$, which gives $\blambda = w_{\infty} = \det(M_d)/ \det(M_{d-1})$.
\end{proof}


\begin{lemm}\label{lem:kernel}
Let $w_{\infty}, s_1,\,\hdots,\,s_d\in \R$ be as given by Lemma~\ref{lem:InftyNodes}. If $y\in \R$ with $F_{\infty}(y)\neq 0$, then there is a quadrature rule for $\mu$ of degree~$2d$ with nodes $y, r_1,\,\hdots,\,r_d\in \R$. Let
\[
q_y \, = \, \prod_{i=1}^{d}(\t-s_i) - \prod_{j=1}^d(\t-r_j)\, \in \, \R[\t]_{\leq\, d-1}.
\]
Then for all $p\in \R[\t]_{\leq\,d-1}$,
\[
\int p \cdot q_y\,\cdot\,(\t-y) \, d\mu \, = \, w_\infty \ev^{d-1}_{\infty}(p).
\]
\end{lemm}


\begin{proof}
By Corollary~\ref{cor:Uniqueness}, there is a quadrature rule for $\mu$ of degree~$2d$ with nodes $y, r_1,\hdots, r_d$ in $\R\cup \{\infty\}$. Since $F_{\infty}(y)\neq 0$ the univariate polynomial $F(\x,y)\in\R[\x]$ has degree $d$ and by the uniqueness of the quadrature rule, $r_1,\hdots, r_d$ are its roots. In particular, $r_j\in \R$. Then for $p\in \R[\t]_{\leq d-1}$, we have that
\begin{align*}
\int p \cdot q_y\,\cdot\,(\t-y) \, d\mu
\ &= \, \int p \cdot \prod_{i=1}^{d}(\t-s_i) \cdot (\t-y) \ d\mu - \int p\,\cdot\, \prod_{j=1}^d(\t-r_j)\,\cdot\,(\t-y) \, d\mu\\
&= \, \int p\,\cdot\,\prod_{i=1}^{d}(\t-s_i) \cdot (\t-y) \ d\mu\\
&= \, w_{\infty}\ev^{2d}_{\infty}\left(p \,\cdot\,\prod_{i=1}^{d}(\t-s_i)\,\cdot\,(\t-y)\right)\\
&= \, w_{\infty}\ev^{d-1}_{\infty}(p).
\end{align*}
The first equality comes from the fact that there is a quadrature rule for $\mu$ of degree~$2d$ with nodes $y,r_1,\,\hdots,\,r_d$. The second follows from the equality in Lemma~\ref{lem:InftyNodes}.
\end{proof}

We now make a special choice of $q$ and use Lemma~\ref{lem:kernel} to greatly simplify $B_{x,y}(p,q)$.

\begin{lemm}\label{lem:q}
Suppose $x,y\in \R$ satisfy $F_{\infty}(x)\neq0$ and $F_{\infty}(y)\neq0$. Take $w_\infty\in \R$ and $q_x, q_y\in \R[\t]_{\leq \,d-1}$ as given by Lemmas~\ref{lem:InftyNodes} and~\ref{lem:kernel} and fix $q = (w_{\infty}, q_x, -q_y)$. Then for any $p = (p_0, p_1, p_2)\in \R\oplus \R[\t]_{\leq\,d-1}\oplus \R[\t]_{\leq\,d-1}$,
\[
B_{x,y}(p,q) \, = \, p_0\cdot \ev_{\infty}(q_x-q_y) +p_0 \cdot (x-y).
\]
\end{lemm}


\begin{proof}
Let $q = (w_{\infty}, q_x, -q_y)$ and suppose $p = (p_0,p_1, p_2) \in \R \oplus \R[\t]_{\leq\,d-1} \oplus \R[\t]_{\leq\,d-1}$. By Lemma~\ref{lem:InftyNodes}, $w_{\infty} $ equals $\det(M_d)/\det(M_{d-1})$. Then, by definition,
\[
B_{x,y}(p,q) =\int \! p_1q_x (x-\t) - p_2q_y (\t-y) d\mu + \ev_{\infty}(w_{\infty}(p_1+p_2)+p_0(q_x-q_y)) +p_0(x-y).\\
\]
Using the properties of $q_x, q_y$ given in Lemma~\ref{lem:kernel}, this simplifies to
\begin{align*}
B_{x,y}(p,q)& = -w_{\infty}\ev_{\infty}(p_1) -w_{\infty} \ev_{\infty}(p_2) +\ev_{\infty}(w_{\infty}(p_1+p_2)\\
& \mkern 240mu +p_0(q_x-q_y)) +p_0(x-y)\\
& = \ev_{\infty}(p_0(q_x-q_y)) +p_0(x-y),
\end{align*}
as claimed.
\end{proof}

\begin{lemm}\label{lem:zard}
Let $H\in\R[\x,\y]$ be a non-zero polynomial with the property that for every $y\in\R$ the polynomial $H(\x,y)\in\R[\x]$ has distinct, real zeros. Then every polynomial vanishing on the real variety of $H$ must be a polynomial multiple of $H$.
\end{lemm}

\begin{proof}
Let $H=H_1\cdots H_r$ where each $H_i$ is irreducible in $\CC[\x,\y]$. By our \linebreak assumption on distinct zeros, the $H_i$ are pairwise coprime. Moreover, each factor $H_i$ belongs to $\R[\x,\y]$. If not, then since $H\in \R[\x,\y]$, both $H_i$ and its complex conjugate $\overline{H_i}$ appear as factors of $H$ and have the same real roots along the line $\y = y$ for $y\in \R$, contradicting the distinctness of these roots. Thus it suffices to show that a polynomial vanishing on the real variety of $H$ must be a polynomial multiple of each $H_i$. Since each $H_i$ satisfies our assumption as well, we can assume without loss of generality that $H$ is irreducible itself.

The real variety of $H$ contains infinitely many points so its Zariski closure in $\CC^2$ is at least one dimensional. By irreducibility of $H$, its real variety Zariski-dense in its complex variety. Now the claim follows from Hilbert's Nullstellensatz.
\end{proof}

Now we are ready to prove Theorem~\ref{thm:main}$\MK$\eqref{list1.4.2}.


\begin{proof}[Proof of Theorem~\ref{thm:main}$\MK$\eqref{list1.4.2}]
Let $F$ be the polynomial given by Theorem~\ref{thm:main}$\MK$\eqref{list1.4.1}. Then $F$ has degree~$2d$ in $\x$ and $\y$, with top degree part $\det(M_{d-1})\x^d\,\y^d$, which is non-zero by the non-degenerateness of $\mu$. Then $(\x-\y)F$ has degree~$2d+1$ with top degree part equal to $\det(M_{d-1})(\x-\y)\x^d\,\y^d$. Also, for every $y\in \R$, the polynomial $(\x-y)F(\x,y)\in \R[\x]$ has distinct, real roots. Lemma~\ref{lem:zard} then implies that any polynomial $G\in \R[\x,\y]$ vanishing on the real variety of $(\x-\y)F$ must be a polynomial multiple of it.

We further claim that the points $(x,y)$ in $V_\R((\x-\y)F)$ with $F_{\infty}(x)\neq0$ and $F_{\infty}(y)\neq0$ are Zariski-dense in $V_\R((\x-\y)F)$. That is, any polynomial $G\in \R[\x,\y]$ vanishing on $V_\R((\x-\y)F)\backslash V_\R(F_{\infty}(\x)F_{\infty}(\y))$ also vanishes on $V_\R((\x-\y)F)$ and is therefore a multiple of $(\x-\y)F$. It suffices to show that $(\x-\y)F$ has no factors in common with $F_{\infty}(\x)F_{\infty}(\y)$. The factors of $F_{\infty}(\x)F_{\infty}(\y)$ are given by Lemma~\ref{lem:InftyNodes}. Suppose for the sake of contradiction that $(\x-s_i)$ divides $(\x-\y)F$ for some $i=1,\hdots, d$. It must be that $(\x-s_i)$ divides $F$. This implies that $F(s_i, s_i)$ equals zero, which contradicts the observation in Remark~\ref{rem:PDdiag} that the matrix $M(s_i,s_i)$ is positive definite and $F(s_i,s_i)= \det(M(s_i,s_i))>0$.


Let $G$ denote the determinant of the $(2d+1)\times (2d+1)$ matrix representing the bilinear form $B_{\x,\y}$. We will show that $(\x-\y)F = c\cdot G$ where $c = (-1)^d \det(M_d)/\linebreak\det(M_{d-1})^2$. Note that $c \cdot G$ is also a polynomial of degree~$2d+1$. Inspection shows that its top degree part is $\det(M_{d-1})(\x-\y)\x^d\y^d$. Thus by the above argument, it suffices to show that $G(x,y)=0$ for all $(x,y)\in V_\R((\x-\y)F)$ with $F_{\infty}(x)F_{\infty}(y)\neq 0$.

Now take $x, y\in \R$ with $(x-y)F(x,y)=0$ and $F_{\infty}(x)F_{\infty}(y)\neq 0$. Let $q_x, q_y\in \R[\t]_{\leq\,d-1}$ be the polynomials given by Lemma~\ref{lem:kernel} and let $q = (w_{\infty}, q_x, -q_y)$. We claim that $q$ is in the kernel of $B_{x,y}$. To see this, let $p = (p_0, p_1, p_2)\in \R\oplus \R[\t]_{\leq\,d-1} \oplus \R[\t]_{\leq\,d-1}$. Then
\[
B_{x,y}(p,q) \, = \, p_0\,\cdot\,\ev_{\infty}(q_x-q_y) +p_0\,\cdot\,(x-y),
\]
by Lemma~\ref{lem:q}. If $x=y$, this is clearly zero. Otherwise $F(x,y)=0$ and by Theorem~\ref{thm:main}$\MK$\eqref{list1.4.1} there exists $r_3, \hdots, r_{d+1}\in \R \cup \{\infty\}$ and a quadrature rule of degree~$2d$ for $\mu$ with nodes $x,y,r_3, \hdots,r_{d+1}$. Since $F_{\infty}(x)\neq 0$, each $r_i\in \R$. Then
\[
q_x \, = \, \prod_{i=1}^{d}(\t-s_i) - (\t-y)\prod_{j=3}^{d+1}(\t-r_j)
\quad \text{and} \quad q_y \, = \, \prod_{i=1}^{d}(\t-s_i) - (\t-x)\prod_{j=3}^{d+1}(\t-r_j).
\]
Expanding and looking at the coefficient of $\t^{d-1}$ reveals that
\[
\ev_{\infty}(q_x) = -\sum_{i=1}^d s_i + y+ \sum_{j=3}^{d+1} r_j
\quad \text{ and } \quad
\ev_{\infty}(q_y) = -\sum_{i=1}^d s_i + x+ \sum_{j=3}^{d+1} r_j.
\]
In particular, $\ev_{\infty}(q_x -q_y) = y-x$, giving $B_{x,y}(p,q)=0$. Since the bilinear form $B_{x,y}$ has a non-zero kernel, the determinant, $G(x,y)$, of its representing matrix is zero.
\end{proof}

\begin{rema}\label{rem:GenEigLinear}
Again this translates the problem of finding the roots of $F(\x,y)$ for fixed $y\in \R$ into a generalized eigenvalue problem. Consider the $(2d+1)\times (2d+1)$ symmetric matrix polynomial
\[
\mathcal{M} \, = \, \mathcal{M}(\x, \y) \, = \,
\begin{pmatrix}
\frac{\det M_{d-1}}{\det M_d} (\x-\y) &e_d^T&e_d^T\\
e_d&\x M_{d-1}-M_{d-1}'&0\\
e_d&0&-\y M_{d-1}+M_{d-1}'
\end{pmatrix}.
\]
Since $M_{d-1}$ is positive definite and $\det(M_{d-1})/\det(M_d)$ is positive, the coefficient of $\x$ in $\mathcal{M}$ is positive semidefinite of rank $d+1$. In particular, for fixed $y\in \R$, we can solve $F(\x,y)=0$ by solving $0=\det(\mathcal{M}(\x,y)) = \det(\x A-B)$, where $A$ is positive semidefinite.
\end{rema}



\begin{exam}[Normal Distribution, $d=3$] \label{ex:Normal2}
Consider again $d=3$ and the normal distribution given in Example~\ref{ex:Normal}. We calculate that $\det(M_3) = 12$ and $\det(M_2) = 2$. The degree~$6$ polynomial $F$ given by Theorem~\ref{thm:main} satisfies $(\x-\y)\linebreak F= -3\det(\mathcal{M})$ where
%{}\small pour la matrix
\[
\mathcal{M} \, = \,
\begin{pmatrix}
\frac{\x-\y}{6} & 0 & 0 & 1& 0 & 0 & 1\\
0& \x & -1 & \x & 0 & 0 & 0 \\
0& -1 & \x & -3 & 0 & 0 & 0 \\
1& \x & -3 & 3 \x & 0 & 0 & 0 \\
0& 0 & 0 & 0 & -\y & 1 & -\y \\
0& 0 & 0 & 0 & 1 & -\y & 3 \\
1& 0 & 0 & 0 & -\y & 3 & -3 \y \\
\end{pmatrix}.
\]
Then $\mathcal{M}(\x,1)$ has the form $\x A - B$ where $A$ is a positive semidefinite matrix of rank four. The determinant $\det(\mathcal{M}(\x,1))$ has four roots, $\x\approx -2.15, -0.52, 1, 2.67$, which are the nodes of a quadrature rule for $\mu$ of degree~$6$. The curve $\det(\mathcal{M})=0$ along with the line $\y=1$ are shown in Figure~\ref{fig:Normal2}.
%\hfill $\diamond$
\end{exam}

\begin{figure}
\begin{center}
%\includegraphics[height=3in]{SexticL}
\includegraphics[height=3in]{Figures/SexticL.pdf} 
\end{center}
\caption{\label{fig:Normal2}
The curve given by $\det(\mathcal{M})=0$ and signs of eigenvalues of the $7\times 7$ matrix $\mathcal{M}$ from Example~\ref{ex:Normal2}. }
\end{figure}


\begin{rema}
Note that the matrix $\mathcal{M}(\x,-\y)$ has the form $\x\,A + \y B + C$ where the matrices $A,B$ are positive semidefinite. For any $a_1, a_2 \in \R_{>0}$ and $b_1,b_2\in \R$, the matrix coefficient of $\t$ in $\mathcal{M}(a_1\t + b_1,-a_2\t + b_2)$ is positive definite, implying that the polynomial $F(a_1\t + b_1,-a_2\t + b_2)\in \R[\t]$ is real-rooted. One can see this in Figures~\ref{fig:Normal} and~\ref{fig:Normal2}, as any line with negative slope intersects the curve $V(F)$ in six real points. It also shows the polynomial $F(\x,-\y)$ is \emph{real stable}~\cite[Proposition~2.4]{wag}. The Helton--Vinnikov theorem~\cite[Theorem~2.2]{hv} then implies that not only $(\x-\y)\,\cdot\,F$ has a $(2d+1)\times(2d+1)$ linear determinantal representation like in the Theorem~\ref{thm:main}$\MK$\eqref{list1.4.2} but also $F$ itself has a determinantal representation $F = \det(\x A + \y B + C)$, where $A,B,C$ are $2d\times2d$ real symmetric matrices and $A$ and $-B$ are positive semidefinite. However it is unclear if there exists such a representation for which the entries of $A,B,C$ are easily calculated from the moments $m_{k}$ of $\mu$, as with $\mathcal{M}$.
\end{rema}

\section{Univariate quadrature rules with more nodes}\label{sec:Generalizations}

It is natural to try to generalize the above discussion to situations where more nodes of a quadrature rule are prescribed. Finding a quadrature rule means specifying an even number of real parameters since each node comes with a weight. We will now consider the following \emph{minimal problem}:

\begin{prob} \label{prob:UnivariateGeneralization}
For integers $n,\ell \geq 1$, we are given $n+2\ell+1$ moments $m_0,\dots, \linebreak m_{n+2\ell}$ of a (positive Borel) measure $\mu$ on the real line that is non-degenerate in degree $n+2\ell$ and $n-1$ real numbers $x_1,\dots, x_{n-1}$. Does there exist a quadrature rule for $\mu$ of degree $n+2\ell$ with $n+\ell$ nodes including $x_1,\dots, x_{n-1}$?
\end{prob}

Specifying the quadrature rule requires $2(\ell+1)+n-1=n+2\ell+1$ parameters, as we have two parameters for each of $\ell+1$ unspecified nodes and one parameter for the weight of each of the $n-1$ specified nodes. Therefore the number of parameters that we have to choose is exactly equal to the number of moments that we need to match.

For $n=1$, the problem is solved by the well-known Gaussian quadrature. The case $n=2$ is the main focus of this paper, and it is classically known that such a measure exists. Unfortunately, for $n=3$ this can fail even if the measure is non-degenerate in every degree. To see this, we need some preparation.


For each integer $0 \leq k \leq n$, consider the quadratic form $p \mapsto \int \t^k \cdot p^2 \ d \mu$ on $\R[\t]_{\leq\,\ell}$, which, with respect to the basis $1,\t,\hdots, \t^{\ell}$, is represented by the $(\ell+1)\times (\ell+1)$ matrix
\[
M^{(k)}_{\ell} \, = \, (m_{i+j+k-2})_{1\,\leq\, i,j\,\leq\,\ell+1}.
\]
Notice that the highest moment of $\mu$ needed to specify these matrices is $m_{n+2\ell}$, achieved by $i = j = \ell+1$ and $k = n$.

For $k\in\N_0$, we denote by $e_k(\x_1,\,\dots,\,\x_n) \in \R[\x_1,\,\hdots,\,\x_n]$ the $k^{\rm th}$ elementary symmetric polynomial in $\x_1,\,\hdots,\,\x_n$, for which the following polynomial identity in $\R[\t]$ holds:
\[
\prod_{i=1}^n(\t-\x_i) \, = \, \sum_{k=0}^n(-1)^ke_k(\x_1,\dots,\x_n)\t^{n-k}.
\]

\begin{prop}\label{multiprop}
Let $\mu$ be a measure on $\R$ that is non-degenerate in degree $n+2\ell\,\geq\,1$ for some $n,\ell \in\N_0$ and let $x_1,\,\hdots,\,x_{n} \in \R$ be distinct. If there is a quadrature rule for $\mu$ of degree $n+2\ell$ with nodes $r_1 = x_1,\,\hdots,\,r_n = x_n, r_{n+1},\,\dots,\,r_{n+\ell}\in\R\cup\{\infty\}$~then
\[
\det\left(\sum_{k=0}^n (-1)^k e_k(x_1,\dots,x_n) M^{(n-k)}_{\ell}\right)=0.
\]
\end{prop}

\begin{proof}
For $X = \{x_1,\,\hdots,\,x_{n}\} \subset \R$, consider the bilinear form $B_{X}$ on $\R[\t]_{\leq\,\ell}$ given by
\[
B_X(p,q) \, = \, \int \, \prod_{k=1}^{n}(\t-x_k) \cdot p\,\cdot\,q \, d\mu
\ \ = \; \sum_{k=0}^n (-1)^k e_k(x_1,\dots,x_n) \cdot \int \t^{n-k}\,\cdot\,p\,\cdot\,q \ d\mu.
\]
With respect to the basis $1,\t,\,\hdots,\, \t^{\ell}$, this is represented by the matrix
\begin{equation}\label{eq:MultiLinearMatrix}
\sum_{k=0}^n (-1)^k e_k(x_1,\,\dots,\,x_n) M^{(n-k)}_{\ell}.
\end{equation}
Suppose that $r_1 = x_1, \hdots, r_n = x_n, r_{n+1},\,\dots,\,r_{n+\ell}\in\R\cup\{\infty\}$ are the nodes of a quadrature rule of degree $n+2\ell$ for $\mu$. As in the proof of Theorem~\ref{thm:main}$\MK$\eqref{list1.4.1}, let $q$ be the unique (up to scaling) nonzero polynomial in $\R[\t]_{\leq\,\ell}$ for which $\ev_{r_i}(q)=0$ for each $i = n+1, \hdots, n+\ell$. Then for every $p\in \R[\t]_{\leq\,\ell}$,
\[
B_X(p,q) \, = \, \int \, \prod_{k=1}^{n}(\t-x_k)\,\cdot\,p\,\cdot\,q \ d\mu \, = \, 0.
\]
Therefore the bilinear form $B_X$ has a nonzero kernel and the determinant of its representing matrix~\eqref{eq:MultiLinearMatrix} is zero.
\end{proof}

Unfortunately, the converse of Proposition~\ref{multiprop} does not hold for $n>2$.

\begin{exam}
Consider the measure given by the exponential distribution on $\R$, whose $k^{\rm th}$ moment is $m_k = k!$ for every $k\in \N_0$. Let $n = 3$, $\ell=3$ so that $n+2\ell=9$. Then in Problem~\ref{prob:UnivariateGeneralization}, we want to build a quadrature rule for $\mu$ of degree~$9$ with at most $ n+\ell=6$ nodes including $n-1=2$ specified nodes $x_1, x_2 \in \R$. However for $x_1=\frac13$ and $x_2=11$, the determinant of the $4\times 4$ matrix in Proposition~\ref{multiprop} equals (up to a positive constant)
\[
137503\x_3^4 -1695024\x_3^3+11282760\x_3^2-41197920\x_3+46998216 \, \in \, \R[\x_3],
\]
which has complex roots $\x_3 \approx 1.87, 5.20, 2.63 \pm i \ 5.31$, only two of which are real. Since $\mu$ is non-degenerate in every degree, any quadrature rule for $\mu$ in degree~$9$ has at $\geq 5$ nodes. However, by Proposition~\ref{multiprop} there are only two possibilities for the remaining $\geq 3$ nodes. Therefore there is no quadrature rule for $\mu$ of degree~$9$ with $\leq 6$ nodes including $x_1=\frac13$ and $x_2=11$.
\end{exam}



